\chapter*{结语\quad 革命的张力}
\addcontentsline{toc}{chapter}{结语\quad 革命的张力} % 手动加入目录
\markboth{结语\quad 革命的张力}{结语\quad 革命的张力} % 设置页眉

循着中共在中央苏区成长、壮大以及遭遇挫折的复杂轨迹一路走来，重构当年这段精彩场景的任务就过程而言已经走到了终点。然而，历史的回廊是如此曲折，个人的认知又是如此有限，因此，过程的完成远远不意味着认知的终结。当我们一路追问，试图去点亮、触碰历史的那些暗角时，常常会不无遗憾地发现，可以照亮的部分是如此有限，而且，就在这照亮的部分中，还有着难以透析的重重光晕。

革命，这个曾经震撼了20世纪中国的“幽灵”，不容分说地攫住了20世纪上半叶众多中国知识者的心灵，革命是20世纪前半期中国历史的主题。从辛亥革命、国民革命到共产革命，一代一代的知识青年感召于国族沦亡危机的呼吁，在屈辱、绝望和不满中奋起，走上革命救国之路。猎猎飘扬的革命大纛展现的是未来、理想和希望，即便其后有暴力、有背叛、有沉滓泛起、有无谓牺牲，仍然掩盖不了革命的光芒，继续革命是革命者期望的克服一切问题的良方。

革命的面貌是如此正义，以致在20世纪的中国，抢夺革命的旗帜成为多种政治势力着力的焦点。最典型的是国民党通过国民革命取得政权，并在中共的共产革命中被视为反革命后，却仍然打着革命的旗帜，以革命党自居。尽管和同时期以苏维埃革命的名义，正在进行社会政治经济关系全面变革的中共相比，完成政权转换使命的国民党事实上已经不再具有现实的革命目标，而且面临着由革命党向执政党转型的实际需要，但革命仍然是其不愿放弃的有号召力的理念。

当革命成为进步和正义的化身时，确实，可以看到许多革命被制造的实例，中国革命也不例外，苏维埃革命时期中共武装对农村的介入就具有一定的空降革命色彩。不过，如果没有20世纪世界革命的大背景，没有席卷全球的民族意识的苏醒，没有中国知识者空前强烈的求变求新愿望，没有技术进步导致的世界政治、经济、观念一体化的趋势，这样的巨大的政治运动可以腾空而起是无法想象的。革命在革命者的手中被创造，革命者的创造历史地生成于当时的现实环境，而革命运动的巨大能量又在不断造就着推展革命的土壤。

对中国革命这样一个如此剧烈地改变、主宰中国命运的政治运动，简单地评判其应该或不应该其实并无实质意义，尤其是在中共领导的共产革命取得政权并对中国社会政治进行了摧枯拉朽式的改造后，更是如此。

革命改变着世界政治版图，改变着国家、社会、家庭、个体，革命尤其是共产革命的张力使这种改变具有为其他政治运动所远远不及的能量，这是革命曾经受到朝圣般欢呼的根由。不过，无论是历史具体情境下的革命事件，还是整体范围内的革命运动，终究还是要受到历史和现实环境的制约。革命是一定情境中人们改变自己命运的利器，但不是包治百病的仙丹，在人类文明进程中，革命甚至也非常态。革命常常以打破现存规则为代价，而规则又是人类社会构成、存续和进步的一个不可或缺的基石，其间所显示的冲突，时时提醒人们在革命令人炫目的张力后面，应有也必有自己的限界。

革命在其所以在，存其所当存。作为人类应慎于使用的天赋权力，革命存在于历史，也存在于当下，存在于革命者的心中，也存在于解读者的脑际。

